% PCT-SOURCE: pdf=187 print=175 kind=heading id=ch4-bibliography
\chapterbibliography
\begingroup
\small
\setlength{\emergencystretch}{2em}
% The source interleaves prose with reference blocks.  Keep the native
% bibliography list for each block while the chapter heading supplies the
% single visible bibliography title.
\providecommand{\bibsection}{}
\renewcommand{\bibsection}{}

% PCT-SOURCE: pdf=187 print=175 kind=prose id=bib-4-intro-local-commutativity
The first proof of the global nature of local commutativity is due to R.~Jost
and O.~Steinmann and appeared in

% PCT-SOURCE: pdf=187 print=175 kind=bibliography id=bib-4-1
\begin{thebibliography}{29}
\setcounter{NAT@ctr}{0}
\bibitem{ch4-1} A.~S.~Wightman, ``Quantum Field Theory and Analytic Functions of
Several Complex Variables,'' \emph{J.~Indian Math.~Soc.}, \textbf{24}, 625
(1960).
\end{thebibliography}

% PCT-SOURCE: pdf=187 print=175 kind=prose id=bib-4-intro-polynomial-algebra
The significance of the polynomial algebra \(\mathcal{P}(\mathcal{O})\)
associated with a region \(\mathcal{O}\) of space-time was first brought out
by R.~Haag.

% PCT-SOURCE: pdf=187 print=175 kind=bibliography id=bib-4-2
\begin{thebibliography}{29}
\setcounter{NAT@ctr}{1}
\bibitem{ch4-2} R.~Haag, ``Discussion des `Axiomes' et des propri\'et\'es
asymptotiques d'une th\'eorie des champs locale avec particules compos\'ees,''
in \emph{Les probl\`emes math\'ematiques de la th\'eorie quantique des
champs}, CNRS, Paris, 1959.
\end{thebibliography}

% PCT-SOURCE: pdf=187 print=175 kind=prose id=bib-4-intro-vacuum-cyclic
The fact that the vacuum is cyclic for \(\mathcal{P}(\mathcal{O})\) was
discovered by Reeh and Schlieder.

% PCT-SOURCE: pdf=187 print=175 kind=bibliography id=bib-4-3
\begin{thebibliography}{29}
\setcounter{NAT@ctr}{2}
\bibitem{ch4-3} H.~Reeh and S.~Schlieder, ``Bemerkungen zur Unit\"ar\"aquivalenz
von Lorentzinvarianten Feldern,'' \emph{Nuovo Cimento}, \textbf{22}, 1051
(1961).
\end{thebibliography}

% PCT-SOURCE: pdf=187 print=175 kind=prose id=bib-4-intro-separating-vectors
The standard theorem on separating vectors for von Neumann algebras is to be
found in

% PCT-SOURCE: pdf=187 print=175 kind=bibliography id=bib-4-4
\begin{thebibliography}{29}
\setcounter{NAT@ctr}{3}
\bibitem{ch4-4} J.~Dixmier, \emph{Les alg\`ebres des op\'erateurs dans l'espace
hilbertien (alg\`ebres de von Neumann)}, Gauthier-Villars, Paris, 1957, p.~6.
\end{thebibliography}

% PCT-SOURCE: pdf=187 print=175 kind=prose id=bib-4-intro-annihilation
The fact that \(\mathcal{P}(\mathcal{O})\) does not contain annihilation
operators was recognized by a number of authors. See Ref.~16 below for
example, or

% PCT-SOURCE: pdf=187 print=175 kind=bibliography id=bib-4-5
\begin{thebibliography}{29}
\setcounter{NAT@ctr}{4}
\bibitem{ch4-5} R.~Jost, ``Properties of Wightman Functions,'' in \emph{Lectures on
Field Theory and the Many-Body Problem}, E.~R.~Caianiello (ed.), Academic
Press, New York, 1961.
\end{thebibliography}

% PCT-SOURCE: pdf=187 print=175 kind=prose id=bib-4-intro-irreducibility
The result of Theorem 4-4, that \(\{E_0,\mathcal{P}(\mathcal{O})\}\) is
irreducible, is contained in Ref.~3. That the irreducibility of the field
operators follows from the cyclicity of the vacuum vector was proved in

% PCT-SOURCE: pdf=187 print=175 kind=bibliography id=bib-4-6-7
\begin{thebibliography}{29}
\setcounter{NAT@ctr}{5}
\bibitem{ch4-6} D.~Ruelle, ``On the Asymptotic Condition in Quantum Field Theory,''
\emph{Helv. Phys. Acta}, \textbf{35}, 147 (1962), Appendix, and also

\bibitem{ch4-7} H.~J.~Borchers, ``On the Structure of the Algebra of Field
Operators,'' \emph{Nuovo Cimento}, \textbf{24}, 214 (1962). The idea of the
proof given here is due to R.~Jost.
\end{thebibliography}

% PCT-SOURCE: pdf=187 print=175 kind=prose id=bib-4-intro-pct
The \emph{PCT} theorem was first stated as : if a relativistic quantum field
theory has space inversion, \(P\), as a symmetry, it must have the product of
charge conjugation and time inversion, \(CT\), as a symmetry also. In this
form, it is proved in

% PCT-SOURCE: pdf=187 print=175 kind=bibliography id=bib-4-8
\begin{thebibliography}{29}
\setcounter{NAT@ctr}{7}
\bibitem{ch4-8} G.~L\"uders, ``On the Equivalence of Invariance under Time Reversal
and under Particle-Anti-Particle Conjugation for Relativistic Field
Theories,'' \emph{Dansk. Mat. Fys. Medd.}, \textbf{28}, 5 (1954).
\end{thebibliography}

% PCT-SOURCE: pdf=187 print=175 kind=prose id=bib-4-intro-pauli
It was Pauli who realized that \emph{PCT} itself is always a symmetry.

% PCT-SOURCE: pdf=187 print=175 kind=bibliography id=bib-4-9
\begin{thebibliography}{29}
\setcounter{NAT@ctr}{8}
\bibitem{ch4-9} W.~Pauli, ``Exclusion Principle, Lorentz Group and Reflection of
Space-Time and Charge,'' p.~30 in \emph{Niels Bohr and the Development of
Physics}, W.~Pauli (ed.), Pergamon Press, New York, 1955.
\end{thebibliography}

% PCT-SOURCE: pdf=188 print=176 kind=prose id=bib-4-intro-jost
The proof given here is that of Jost:

% PCT-SOURCE: pdf=188 print=176 kind=bibliography id=bib-4-10-11
\begin{thebibliography}{29}
\setcounter{NAT@ctr}{9}
\bibitem{ch4-10} R.~Jost, ``Eine Bemerkung zum CTP Theorem,'' \emph{Helv. Phys.
Acta}, \textbf{30}, 409 (1957). This paper was the starting point for many
of the applications given in the present chapter. The connection between
\emph{PCT} invariance and weak local commutativity introduced here is further
discussed in

\bibitem{ch4-11} F.~J.~Dyson, ``On the Connection of Weak Local Commutativity and
Regularity of Wightman Functions,'' \emph{Phys. Rev.}, \textbf{110}, 579
(1958).
\end{thebibliography}

% PCT-SOURCE: pdf=188 print=176 kind=prose id=bib-4-intro-pauli-pct
It is interesting that in his paper of 1940, on spin and statistics, Pauli
actually proves what can be regarded as ``classical'' \emph{PCT} invariance
for the free field theories he considers: the substitution rule (1-51) is
shown to leave the equations invariant if those equations are invariant under
the restricted Lorentz group and are linear. The remaining feature of the
substitution rule for \emph{PCT} in quantum mechanics, the reversal of order,
first appears in the following paper of Schwinger:

% PCT-SOURCE: pdf=188 print=176 kind=bibliography id=bib-4-12
\begin{thebibliography}{29}
\setcounter{NAT@ctr}{11}
\bibitem{ch4-12} J.~Schwinger, ``On the Theory of Quantized Fields I,''
\emph{Phys. Rev.}, \textbf{82}, 914 (1951). However, readers of this paper
did not generally recognize that it stated or proved the \emph{PCT} theorem.
In the free field case considered by Pauli no operator products occur so the
``classical'' \emph{PCT} invariance he proved for the equations is the same as
the full quantum mechanical \emph{PCT} invariance.
\end{thebibliography}

% PCT-SOURCE: pdf=188 print=176 kind=prose id=bib-4-intro-spin-statistics
The spin-statistics theorem has a long history. In a fairly general form
(general spin, but for free fields) it appears in

% PCT-SOURCE: pdf=188 print=176 kind=bibliography id=bib-4-13-14
\begin{thebibliography}{29}
\setcounter{NAT@ctr}{12}
\bibitem{ch4-13} M.~Fierz, ``\"Uber die relativistische Theorie kr\"aftefreier
Teilchen mit beliebigem Spin,'' \emph{Helv. Phys. Acta}, \textbf{12}, 3
(1939).

\bibitem{ch4-14} W.~Pauli, ``On the Connection between Spin and Statistics,''
\emph{Phys. Rev.}, \textbf{58}, 716 (1940).
\end{thebibliography}

% PCT-SOURCE: pdf=188 print=176 kind=prose id=bib-4-intro-general-proofs
The first proofs using only the general assumptions of Chapter 3 are due to

% PCT-SOURCE: pdf=188 print=176 kind=bibliography id=bib-4-15-16
\begin{thebibliography}{29}
\setcounter{NAT@ctr}{14}
\bibitem{ch4-15} G.~L\"uders and B.~Zumino, ``Connection between Spin and
Statistics,'' \emph{Phys. Rev.}, \textbf{110}, 1450 (1958); and

\bibitem{ch4-16} N.~Burgoyne, ``On the Connection of Spin with Statistics,''
\emph{Nuovo Cimento}, \textbf{8}, 807 (1958). Our treatment here follows the
latter.
\end{thebibliography}

% PCT-SOURCE: pdf=188 print=176 kind=prose id=bib-4-intro-dellantonio
The significant point covered by Theorem 4-8 was first made clear by
Dell'Antonio:

% PCT-SOURCE: pdf=188 print=176 kind=bibliography id=bib-4-17
\begin{thebibliography}{29}
\setcounter{NAT@ctr}{16}
\bibitem{ch4-17} G.~F.~Dell'Antonio, ``On the Connection of Spin with Statistics,''
\emph{Ann. Phys.}, \textbf{16}, 153 (1961).
\end{thebibliography}

% PCT-SOURCE: pdf=188 print=176 kind=prose id=bib-4-intro-fields
The commutation relations between different fields were discussed in the old
framework by a number of authors of whom we quote only

% PCT-SOURCE: pdf=188 print=176 kind=bibliography id=bib-4-18
\begin{thebibliography}{29}
\setcounter{NAT@ctr}{17}
\bibitem{ch4-18} G.~L\"uders, ``Vertauschungs\-relationen zwischen verschiedenen
Feldern,'' \emph{Z. Naturforsch.}, \textbf{13a}, 254 (1958). L\"uders was the
first to use systematically the vector space \(V\) described in the text. The
notion of Klein transformation which is used by these authors was
systematically used by O.~Klein in another context:
\end{thebibliography}

% PCT-SOURCE: pdf=189 print=177 kind=bibliography id=bib-4-19
\begin{thebibliography}{29}
\setcounter{NAT@ctr}{18}
\bibitem{ch4-19} O.~Klein, ``Quelques remarques sur le traitement approximatif du
probl\`eme des \'electrons dans un r\'eseau cristallin par la m\'ecanique
quantique,'' \emph{J. Phys. Radium}, \textbf{9}, 1 (1938).
\end{thebibliography}

% PCT-SOURCE: pdf=189 print=177 kind=prose id=bib-4-intro-jordan-wigner
It also appears in

% PCT-SOURCE: pdf=189 print=177 kind=bibliography id=bib-4-19a
\begin{thebibliography}{29}
\setcounter{NAT@ctr}{19}
\item[19a.] P.~Jordan and E.~Wigner, ``\"Uber das Paulische
\"Aquivalenzverbot,'' \emph{Z. Physik.}, \textbf{47}, 631 (1928).
\end{thebibliography}

% PCT-SOURCE: pdf=189 print=177 kind=prose id=bib-4-intro-araki
The treatment of the commutation relations between different fields followed
here is due to Araki.

% PCT-SOURCE: pdf=189 print=177 kind=bibliography id=bib-4-20
\begin{thebibliography}{29}
\setcounter{NAT@ctr}{19}
\bibitem{ch4-20} H.~Araki, ``Connection of Spin with Commutation Relations,''
\emph{J. Math. Phys.}, \textbf{2}, 267 (1961).
\end{thebibliography}

% PCT-SOURCE: pdf=189 print=177 kind=prose id=bib-4-intro-haag
Haag's theorem originates in

% PCT-SOURCE: pdf=189 print=177 kind=bibliography id=bib-4-21
\begin{thebibliography}{29}
\setcounter{NAT@ctr}{20}
\bibitem{ch4-21} R.~Haag, ``On Quantum Field Theory,'' \emph{Dan. Mat. Fys.
Medd.}, \textbf{29}, 12 (1955).
\end{thebibliography}

% PCT-SOURCE: pdf=189 print=177 kind=prose id=bib-4-intro-generalized-haag
The generalized Haag's theorem was first proved in

% PCT-SOURCE: pdf=189 print=177 kind=bibliography id=bib-4-22
\begin{thebibliography}{29}
\setcounter{NAT@ctr}{21}
\bibitem{ch4-22} D.~Hall and A.~S.~Wightman, ``A Theorem on Invariant Analytic
Functions with Applications to Relativistic Quantum Field Theory,''
\emph{Mat. Fys. Medd. Dan. Vid. Selsk.}, \textbf{31}, 5 (1957).
\end{thebibliography}

% PCT-SOURCE: pdf=189 print=177 kind=prose id=bib-4-intro-jost-schroer
The simple proof of Theorem 4-15 given here is due to Jost and Schroer
(Ref.~5). A different treatment is found in

% PCT-SOURCE: pdf=189 print=177 kind=bibliography id=bib-4-23
\begin{thebibliography}{29}
\setcounter{NAT@ctr}{22}
\bibitem{ch4-23} O.~W.~Greenberg, ``Haag's Theorem and Clothed Operators,''
\emph{Phys. Rev.}, \textbf{115}, 706 (1959).
\end{thebibliography}

% PCT-SOURCE: pdf=189 print=177 kind=prose id=bib-4-intro-federbush
A proof of Haag's theorem similar to that of Jost and Schroer was given by
Federbush and Johnson:

% PCT-SOURCE: pdf=189 print=177 kind=bibliography id=bib-4-24
\begin{thebibliography}{29}
\setcounter{NAT@ctr}{23}
\bibitem{ch4-24} P.~G.~Federbush and K.~A.~Johnson, ``The Uniqueness of the
Two-Point Function,'' \emph{Phys. Rev.}, \textbf{120}, 1926 (1960).
\end{thebibliography}

% PCT-SOURCE: pdf=189 print=177 kind=prose id=bib-4-intro-borchers
Equivalence classes of local fields were first described in

% PCT-SOURCE: pdf=189 print=177 kind=bibliography id=bib-4-25
\begin{thebibliography}{29}
\setcounter{NAT@ctr}{24}
\bibitem{ch4-25} H.~J.~Borchers, ``\"Uber die Mannig\-faltigkeit der inter\-polierenden
Felder zu einer kausalen S-Matrix,'' \emph{Nuovo Cimento}, \textbf{15}, 784
(1960).
\end{thebibliography}

% PCT-SOURCE: pdf=189 print=177 kind=prose id=bib-4-intro-equivalence-theorems
Examples of relatively local fields which have the same S-matrix when
evaluated in perturbation theory were found by Chisholm, Salam, and
Kamefuchi. References to their papers can be found in

% PCT-SOURCE: pdf=189 print=177 kind=bibliography id=bib-4-26
\begin{thebibliography}{29}
\setcounter{NAT@ctr}{25}
\bibitem{ch4-26} S.~Kamefuchi, L.~O'Raifeartaigh, and Abdus Salam, ``Change of
Variables and Equivalence Theorems in Quantum Field Theories,'' \emph{Nucl.
Phys.}, \textbf{28}, 529 (1961).
\end{thebibliography}

% PCT-SOURCE: pdf=189 print=177 kind=prose id=bib-4-intro-free-hermitian
The equivalence class of a free hermitian scalar field was determined
independently by B.~Schroer (unpublished) and H.~Epstein.

% PCT-SOURCE: pdf=189 print=177 kind=bibliography id=bib-4-27
\begin{thebibliography}{29}
\setcounter{NAT@ctr}{26}
\bibitem{ch4-27} H.~Epstein, ``On the Borchers Class of a Free Field,''
\emph{Nuovo Cimento}, \textbf{27}, 886 (1963).
\end{thebibliography}

% PCT-SOURCE: pdf=189 print=177 kind=prose id=bib-4-intro-statistics
The literature on statistics different from Fermi or Bose statistics can be
traced from the references in

% PCT-SOURCE: pdf=189 print=177 kind=bibliography id=bib-4-28
\begin{thebibliography}{29}
\setcounter{NAT@ctr}{27}
\bibitem{ch4-28} O.~W.~Greenberg and A.~Messiah, ``Are There Particles in Nature
Other Than Bosons or Fermions?'' \emph{Phys. Rev.}, \textbf{136}, B248 (1964).
\end{thebibliography}

% PCT-SOURCE: pdf=190 print=178 kind=prose id=bib-4-intro-haag-ruelle
For a systematic account of the Haag-Ruelle theory, and for further study,
see

% PCT-SOURCE: pdf=190 print=178 kind=bibliography id=bib-4-29
\begin{thebibliography}{29}
\setcounter{NAT@ctr}{28}
\bibitem{ch4-29} R.~Jost, \emph{General Theory of Quantized Fields (Lectures in
Applied Mathematics IV: Proceedings of the Summer Seminar, Boulder, Colorado,
1960)}, American Math. Soc., Providence, R.I., 1965.
\end{thebibliography}
\endgroup
